\chapter{Cosmetic Procedures}

\section*{Definition and Overview}
Cosmetic procedures are treatments performed to improve a person's appearance. They range from non-invasive treatments such as dermal fillers, Botox, and laser therapy to surgical operations such as breast augmentation, liposuction, and facelifts. Cosmetic procedures differ from reconstructive surgery, which restores function or corrects disfigurement. Some cosmetic treatments are performed by doctors and plastic surgeons, while others, such as injectable fillers, may be offered in non-medical settings; in Australia many of these are regulated as medical procedures. Anyone considering a cosmetic procedure should understand what it involves, its risks, the qualifications of the practitioner, and the realistic outcome.

\section*{Epidemiology}
Cosmetic procedures are very common worldwide, and demand has grown steadily. In Australia, the most popular treatments include injectables such as Botox and dermal fillers, laser hair removal, and surgical procedures such as breast augmentation, liposuction, and rhinoplasty. The majority of patients are women, but the number of men seeking cosmetic treatments is rising. Social media and the availability of cheaper treatments have increased demand, particularly among younger adults, and there has been a parallel rise in complications, often linked to unqualified practitioners and unregulated products.

\section*{Aetiology and Pathophysiology}
People seek cosmetic procedures for many reasons, and the treatments work through different mechanisms.

\begin{itemize}
    \item \textbf{Ageing changes:} loss of collagen, fat, and skin elasticity with age, treated by fillers, Botox, and facelift surgery
    \item \textbf{Body concerns:} dissatisfaction with breast size, body contours, or facial features
    \item \textbf{Injectable treatments:} Botox relaxes muscles to smooth wrinkles; fillers restore volume
    \item \textbf{Laser and light treatments:} target pigmentation, redness, hair, and skin texture
    \item \textbf{Surgical procedures:} remove excess skin and fat, reshape tissues, or place implants
    \item \textbf{Psychological factors:} low self-esteem, body image concerns, and pressure from media or partners
    \item \textbf{Body dysmorphic disorder:} a genuine mental health condition in which cosmetic procedures rarely help and can worsen distress
\end{itemize}

\section*{Clinical Features}
Rather than symptoms, cosmetic consultations are driven by a person's appearance concerns and expectations.

\begin{itemize}
    \item Dissatisfaction with a specific facial or body feature
    \item Visible ageing changes such as wrinkles, volume loss, or sagging skin
    \item Unwanted body fat or lax skin after weight loss or pregnancy
    \item Concerns about scars, pigmentation, or hair
    \item A desire to match an appearance seen in media or social media
    \item Ongoing dissatisfaction despite previous procedures
\end{itemize}

\section*{Differential Diagnosis}
A responsible consultation identifies who is and is not a suitable candidate.

\begin{itemize}
    \item Realistic cosmetic goals versus expectations a procedure cannot meet
    \item Ageing changes treatable by non-surgical means versus those needing surgery
    \item Body image concerns within a normal range versus body dysmorphic disorder
    \item Medical conditions that increase risk, such as poorly controlled diabetes or bleeding disorders
    \item Signs of depression, anxiety, or low self-esteem that need psychological support
    \item Unlicensed, counterfeit, or unsuitable products versus regulated medical devices
\end{itemize}

\section*{When to Seek Medical Care}
Before any cosmetic procedure, have a medical assessment and check the practitioner's qualifications and the clinic's credentials. Seek immediate care for complications such as infection, allergic reaction, or skin necrosis.

\textbf{Call 000 (or local emergency services) for:}
\begin{itemize}
    \item Sudden severe pain, swelling, or difficulty breathing after a procedure
    \item Signs of a serious allergic reaction: swelling of the face or throat, wheezing, or collapse
    \item A spreading skin infection or fever after surgery
    \item Chest pain or shortness of breath, which may indicate a blood clot after surgery
\end{itemize}

\section*{Diagnosis and Investigations}
Assessment before a cosmetic procedure focuses on safety, suitability, and consent.

\begin{itemize}
    \item \textbf{Medical history:} general health, medicines, allergies, previous reactions, and bleeding risk
    \item \textbf{Examination:} assessment of skin quality, anatomy, and the features to be treated
    \item \textbf{Photography:} standardised images for planning and comparison
    \item \textbf{Psychological screening:} discussion of expectations and mental health, with referral where appropriate
    \item \textbf{Product and practitioner verification:} confirmation that the practitioner is qualified and products are approved and genuine
    \item \textbf{Informed consent:} a clear explanation of risks, benefits, costs, and expected results
\end{itemize}

\section*{Management}
The appropriate treatment depends on the concern, the person's anatomy, and their goals.

\begin{itemize}
    \item \textbf{Injectable treatments:} Botox and dermal fillers for wrinkles and volume loss
    \item \textbf{Skin treatments:} lasers, peels, and microneedling for pigmentation, texture, and hair
    \item \textbf{Non-surgical body treatments:} cryolipolysis and other devices for localised fat
    \item \textbf{Minor surgical procedures:} eyelid surgery, liposuction, and other day procedures
    \item \textbf{Major surgery:} breast augmentation, facelift, and abdominoplasty under anaesthetic
    \item \textbf{Aftercare:} wound care, sun protection, and follow-up visits
    \item \textbf{Psychological support:} counselling for body image concerns where appropriate
    \item \textbf{Revision:} further procedures if results are unsatisfactory, with realistic discussion
\end{itemize}

\section*{Complications}
\begin{itemize}
    \item Bruising, swelling, pain, and infection
    \item Allergic reactions to products, anaesthetic, or dressings
    \item Scarring, skin necrosis, or pigment changes
    \item Asymmetry, lumpiness, or an unnatural result
    \item Nerve injury causing numbness, weakness, or paralysis
    \item Complications of anaesthesia and surgery, including bleeding and blood clots
    \item Psychological distress, including dissatisfaction and worsening body image
    \item Complications from unregulated products or unqualified practitioners, including permanent harm
\end{itemize}

\section*{Prognosis and Outcomes}
Most cosmetic procedures achieve their intended improvement when performed by qualified practitioners on suitable patients. Non-invasive treatments give temporary results lasting months, while surgical results are longer lasting but carry greater risk. Satisfaction depends strongly on realistic expectations and honest communication. People with body dysmorphic disorder are rarely helped and may become more distressed, which is why responsible practitioners screen for it. Choosing a qualified, registered practitioner and an accredited facility is the single most important step in achieving a good outcome.

\section*{Prevention and Self-Care}
\begin{itemize}
    \item Choose a registered medical practitioner with specific cosmetic training
    \item Verify that injectable products are genuine, approved, and used from sealed containers
    \item Have a full consultation, including risks, costs, and recovery, before agreeing
    \item Ask to see before-and-after photos and the practitioner's credentials
    \item Be wary of deals, pressure sales, and treatments offered in non-medical settings
    \item Consider whether your expectations are realistic, and seek counselling if body image concerns are significant
    \item Follow aftercare instructions, including sun protection and wound care
    \item Do not travel abroad for procedures that cannot be safely followed up
\end{itemize}

\section*{Sources and Further Reading}
The following reputable sources provide further information for healthcare students and patients seeking reliable, current advice about cosmetic procedures and their safety.

\begin{itemize}
    \item Australian Health Practitioner Regulation Agency (AHPRA). \textit{Regulation of cosmetic procedures.} https://www.ahpra.gov.au/
    \item Therapeutic Goods Administration. \textit{Injectable cosmetic products: safety information.} https://www.tga.gov.au/
    \item NHS. \textit{Cosmetic procedures: what to consider.} https://www.nhs.uk/conditions/cosmetic-procedures/
    \item Royal Australasian College of Surgeons. \textit{Cosmetic surgery patient information.} https://www.surgeons.org/
    \item American Society of Plastic Surgeons. \textit{Cosmetic procedure information.} https://www.plasticsurgery.org/
    \item Healthdirect Australia. \textit{Cosmetic procedures and surgery.} https://www.healthdirect.gov.au/
\end{itemize}
